%!TEX root = main.tex
\section{Introduction}
\label{sec:introduction}

\PARstart{W}{earable} biomedical devices---continuous glucose monitors, cardiac rhythm analyzers, inertial fall predictors---generate vast amounts of physiological data transmitted wirelessly to cloud platforms~\cite{wearable-health}. Regulatory frameworks (HIPAA, GDPR) impose strict requirements on data provenance, yet current device architectures lack mechanisms for the device itself to \emph{prove} that a measurement is authentic without revealing raw signals. Zero-Knowledge Proofs (ZKPs)~\cite{groth16,plonk} resolve this tension: a prover (the sensor) convinces a verifier (a server) that a statement holds---e.g., that a glucose trajectory was computed correctly---without disclosing the underlying data. Elliptic Curve Cryptography (ECC)~\cite{ecc-iot} underpins the most compact constructions, keeping proof sizes small for bandwidth-constrained links.

Deploying ZKP on hardware demands \emph{reconfigurability}. Different proof systems (Groth16~\cite{groth16}, PLONK~\cite{plonk}, STARKs~\cite{stark-anatomy}) use different prime fields (31-bit BabyBear, 64-bit Goldilocks, 256-bit BN254), S-box exponents ($d\!\in\!\{3,5,7\}$), and elliptic curves. Plonky3~\cite{babybear} introduced BabyBear only in 2024, replacing the Goldilocks field of Plonky2. An ASIC designed for one configuration becomes obsolete when the next proof system changes parameters. Coarse-Grained Reconfigurable Arrays (CGRAs) occupy the right design point: unlike FPGAs, they reconfigure at the instruction level with microsecond latency; unlike ASICs, they remain fully reprogrammable~\cite{cgra-survey}. In our research, we target the X-HEEP framework~\cite{xheep}, an extensible
RISC-V based SoC designed for ultra low power edge computing, capable of interfacing with custom peripherals like accelerators. In particular, applications of the X-HEEP framework, such as  HEEPocrates and HEEPsilon feature a CGRA integrated via the Extendible Accelerator Interface (XAIF): the OpenEdgeCGRA~\cite{openedge-cgra},
a $4 \times 4$ PE mesh with a 32-bit datapath,
which we consider in this paper.

The computational bottleneck in ZKP is the hash function. We focus our research on the Poseidon2 framework~\cite{poseidon2}, which implement a hash that minimizes finite-field multiplications per permutation making it highly suitable for our data privacy use case (a biomedical wereable). In any case, its S-box layer---$x^d \Mod p$---accounts for over 90\% of computation on a CGRA with 32-bit datapath. For 31-bit BabyBear ($\babybear$, $d\!=\!7$), each S-box requires two modular squarings and two modular multiplications, each involving Barrett reduction~\cite{barrett1986} with 62-bit intermediates decomposed across 32-bit PEs. Yet programming the CGRA is the real bottleneck: the developer must specify, for each of the $N\times M$~PEs and each clock cycle, the exact instruction, registers, and neighbor links. For the S-box layer this spans \lines{\MetricLocManual}, took \MetricSboxKernelManualAuthoringTime{} of expert effort, and a single misplaced operand produces silent errors detectable only through full simulation. Automatic compilers~\cite{compigra,wang2025mlir} avoid this burden but cannot express cross-PE carry chains, toroidal routing, or branchless conditional selects that are essential for modular arithmetic.

In this paper we present \castm{} (CGRA API \& Spatial Temporal Mapper),
which occupies the middle ground: the programmer retains explicit control over \emph{which PE} executes \emph{which instruction} at \emph{which cycle}, but through parametric functions, loops, labeled subroutines, and architecture-aware operators instead of raw CSV encodings. A four-stage deterministic pipeline
produces simulator-native CSV with structured error diagnostics. More precisely, we make the following contributions: \textbf{(C1)} We introduce \castm{}, the first API targeting spatial-temporal CGRA programming for cryptographic workloads (Section~\ref{sec:openedge}).
\textbf{(C2)} We map S-box layer ($x^7 \Mod p$, BabyBear) onto OpenEdgeCGRA and validate the resulting flow against independent C/C++ references (Section~\ref{subsec:sbox-openedge}). \textbf{(C3)} We show \castm{} reduces source code by approximately \timex{\MetricLocCompressionVsManual} vs.\ manual CSV while achieving \pct{\MetricSboxCastmCycleReductionPct} fewer clock cycles on the S-box kernel (Section~\ref{subsec:sbox-comparison}). \textbf{(C4)} We identify multi-batch configuration reload as the dominant bottleneck and evaluate increasing Configuration Register File --CRF-- capacity to enable full Poseidon2 deployment on CGRA (Section~\ref{subsec:creg-analysis}). \textbf{(C5)} We show that throughput and energy favor different deployments of Poseidon2 on the CPU and the CGRA (Section~\ref{subsec:rtl-poseidon2}).
